\documentclass[dvisvgm,tikz,border=3pt]{standalone}
\usepackage{amsmath,amssymb}
\usepackage{pgfplots}
\usepackage{CJKutf8}
\pgfplotsset{compat=1.18}
\usetikzlibrary{arrows.meta,calc,positioning}

\definecolor{themebg}{HTML}{2e362c}
\definecolor{themefg}{HTML}{edf4e8}
\definecolor{themegray}{HTML}{9ea897}
\definecolor{themecurve}{HTML}{7eb6ff}
\definecolor{themealert}{HTML}{ff7b7b}
\definecolor{themeamber}{HTML}{f5b942}

\pgfdeclarelayer{background}
\pgfdeclarelayer{foreground}
\pgfsetlayers{background,main,foreground}

\begin{document}
\begin{CJK*}{UTF8}{gbsn}
\pagecolor{themebg}
\begin{tikzpicture}[color=themefg, text=themefg]

\begin{axis}[
    width=12.2cm,
    height=9.0cm,
    axis lines=middle,
    axis line style={color=themefg, -{Stealth}},
    tick style={color=themefg},
    tick label style={color=themefg, font=\scriptsize},
    every axis label/.style={color=themefg},
    xmin=-0.5, xmax=4.8,
    ymin=-0.5, ymax=4.6,
    xtick={1,2,3,4},
    ytick={1,2,3,4},
    clip=false,
    xlabel={$x$},
    ylabel={$y$},
    title style={font=\footnotesize\bfseries, at={(0.5,1.06)}, anchor=south},
    title={Lagrange 乘数法几何本质：目标等高线与约束曲线在条件极值点必然相切},
]

% 1. 目标函数等高线族 f(x,y) = c
\draw[themecurve!50!themebg, line width=1.0pt] (axis cs:1,1) circle [radius=0.7cm];
\node[text=themecurve!70!themebg, font=\tiny, fill=themebg, inner sep=0.5pt] at (axis cs:1.8,1.4) {$f=c_1$};

\draw[themecurve!70!themebg, line width=1.2pt] (axis cs:1,1) circle [radius=1.3cm];
\node[text=themecurve!85!themebg, font=\tiny, fill=themebg, inner sep=0.5pt] at (axis cs:2.4,1.8) {$f=c_2$};

\draw[themecurve, line width=2.0pt] (axis cs:1,1) circle [radius=1.9cm];
\node[text=themecurve, font=\scriptsize\bfseries, fill=themebg, inner sep=0.8pt] at (axis cs:2.9,1.05) {$f=c^*$ (极值等高线)};

\draw[themecurve!50!themebg, line width=1.0pt] (axis cs:1,1) circle [radius=2.5cm];
\node[text=themecurve!70!themebg, font=\tiny, fill=themebg, inner sep=0.5pt] at (axis cs:3.6,2.8) {$f=c_4$};

% 2. 约束曲线 g(x,y) = 0 (切于 P*)
\addplot[themeamber, line width=2.2pt, domain=0.8:3.9, samples=50] {4.68 - x};
\node[text=themeamber, font=\small\bfseries, fill=themebg, inner sep=1pt, anchor=south west] at (axis cs:3.4,1.4) {约束曲线 $g(x,y)=0$};

\begin{pgfonlayer}{foreground}
    % 条件极值点 P*
    \coordinate (Pstar) at (axis cs:2.34, 2.34);
    \addplot[only marks, mark=*, mark size=3pt, fill=themealert, draw=themefg] coordinates {(2.34, 2.34)};
    \node[text=themealert, font=\small\bfseries, fill=themebg, inner sep=1pt, anchor=south east] at (axis cs:2.25, 2.45) {条件极值点 $P^*$};

    % 梯度向量 grad f (垂直于等高线，沿径向)
    \draw[-{Stealth[length=3mm]}, themecurve, line width=2pt] (Pstar) -- ++(axis cs:1.0, 1.0) node[anchor=south west, font=\scriptsize\bfseries] {$\nabla f(P^*)$};

    % 梯度向量 grad g (垂直于约束线，法向)
    \draw[-{Stealth[length=3mm]}, themeamber, line width=2pt] (Pstar) -- ++(axis cs:0.65, 0.65) node[anchor=north west, font=\scriptsize\bfseries] {$\nabla g(P^*)$};

    % 相交非极值点 A, B
    \addplot[only marks, mark=o, mark size=2.5pt, line width=1.2pt, themegray] coordinates {(1.3, 3.38) (3.38, 1.3)};
    \node[text=themegray, font=\tiny, fill=themebg, inner sep=0.6pt, anchor=south] at (axis cs:1.3, 3.5) {穿越点 (非极值)};

    % 结论框放置在左上，偏右避开 y 轴
    \node[text=themefg, font=\scriptsize, fill=themebg, inner sep=1.5pt, anchor=north west] at (axis cs:0.2, 4.35) {
        \begin{tabular}{l}
        \textbf{相切条件}：等高线切线 $\parallel$ 约束切线\\
        $\implies$ 法向量共线：$\boxed{\nabla f(P^*) = \lambda \nabla g(P^*)}$\\
        若不相切（穿越），则沿约束仍可单调升降
        \end{tabular}
    };
\end{pgfonlayer}

\end{axis}
\end{tikzpicture}
\end{CJK*}
\end{document}
